\subsection{CM nonvanishing with a rational principal-lattice bound}
\label{cmnonzero:sec:nonvanishing}
The possible zero denominators in the product and genus formulas can be
excluded uniformly. Non-elliptic CM nonvanishing is classical, through
Kohnen's conjugation argument~\cite{cmnonzero:Kohnen}; see the introduction
of Chiriac--Jorza~\cite{cmnonzero:ChiriacJorza} for its place in zero theory.
The following proof supplies an explicit principal-point bound, certified
using rational arithmetic. It applies to nonmaximal orders as well.

\begin{theorem}[Nonvanishing at every non-elliptic CM point]
\label{cmnonzero:thm:allCM}
Let $D<-4$ be any quadratic-order discriminant. For every CM point $\tau$
of discriminant $D$ and every even integer $w\ge4$,
\[
 G_w(\tau)\ne0.
\]
For the principal reduced point
\[
 \tau_D=\begin{cases}i\sqrt{|D|}/2,&D\equiv0\pmod4,\\
 (-1+i\sqrt{|D|})/2,&D\equiv1\pmod4,\end{cases}
\]
one has the weight- and discriminant-uniform bound
$|G_w(\tau_D)|>7/20$.
\end{theorem}

\begin{proof}
The principal lattice has, for $n\ne0$, squared lengths bounded below by
one of the two integer quadratic forms
\[
 q_7(m,n)=m^2-mn+2n^2,\qquad q_8(m,n)=m^2+2n^2.
\]
Indeed the odd principal discriminant has $|D|\ge7$, and the even one
has $|D|\ge8$. Both forms are at least two when $n\ne0$. Hence for every
even $w\ge4$ the off-axis absolute sum is bounded by
\[
 T_d=\sum_{n\ne0}\sum_{m\in\mathbb Z}q_d(m,n)^{-2},\qquad d=7,8.
\]
We prove $T_7<33/20$ and $T_8<11/8$ by finite rational sums and explicit
infinite-tail bounds.

For any $A>0$ and real shift $s$, layer-cake integration gives
\[
 \sum_{m\in\mathbb Z}\frac1{((m-s)^2+A^2)^2}
 \le A^{-4}+\frac{\pi}{2A^3}.
\]
To see this, each superlevel set of the even decreasing function is an
interval. Its number of translated lattice points is at most its length
plus one. Integrating that inequality bounds the lattice sum by the
integral of the function plus its maximum. The integral is $\pi/(2A^3)$.

Use $N=10$, $M=20$ in the finite rectangle. The imaginary coordinate in
both lattices is greater than one. Thus the part with $|n|>N$, summed
over every $m$, is less than
\[
 \frac2{3N^3}+\frac2{N^2},
\]
by $\pi<4$, $\sum_{n>N}n^{-4}\le1/(3N^3)$ and
$\sum_{n>N}n^{-3}\le1/(2N^2)$. For $1\le |n|\le N$ and $|m|>M$,
the squared length is at least $(|m|-N/2)^2$. Integral comparison bounds
this part by $4N/[3(M-N/2)^3]$. The finite certificate therefore verifies
\[
 2\sum_{n=1}^{10}\sum_{m=-20}^{20}q_7(m,n)^{-2}
       +\frac{31}{1500}+\frac8{2025}<\frac{33}{20},
\]
and the same expression with $q_8$ is less than $11/8$. Every term and
comparison in these two finite inequalities is rational; the full fractions
are recorded in \texttt{verification/CM-nonvanishing.json}.

The $n=0$ axis contributes exactly $2\zeta(w)>2$. The triangle inequality
now gives $|G_w(\tau_D)|>2-33/20=7/20$ in the odd case, and the stronger
$5/8$ in the even case.

Suppose a CM value at the same discriminant vanished. By
\eqref{cmnorm:eq:Fw}, its singular modulus would be a zero of the rational
polynomial $F_w$. The class polynomial $H_D$ is the minimal polynomial of
that singular modulus, including for a nonmaximal order
\cite{cmnorm:Sutherland,cm:Cox}. Thus $H_D$ divides $F_w$, and the
principal singular modulus is also a zero. Since $\Delta$ never vanishes
on the upper half-plane, this would make $G_w(\tau_D)=0$, contradicting
the bound. Modular transformations multiply a positive-weight value by a
nonzero factor, so the conclusion covers every representative.
\end{proof}

Consequently all class and genus resultants in the following formulas are
nonzero at $D<-4$. Their positive roots and quotients are defined at every
even weight. The excluded discriminants $-3,-4$ retain their familiar
elliptic vanishing patterns.
